
\section{Threat Model}


\begin{figure}[H]
\centering
\begin{tikzpicture}[every node/.style={figlabel}]
\path[figpanel] (-5.3,-3.6) rectangle (5.3,3.75);
\node[figtitle, anchor=west] at (-4.85,3.25) {TRUST BOUNDARY};
\fill[BitRose!5] (0,0) circle (2.95cm);
\draw[BitRose!35,line width=0.9pt] (0,0) circle (2.95cm);
\fill[BitGold!8] (0,0) circle (2.0cm);
\draw[BitGold!55!black,line width=0.9pt] (0,0) circle (2.0cm);
\fill[BitTeal!10] (0,0) circle (1.05cm);
\draw[BitTeal!72,line width=0.95pt] (0,0) circle (1.05cm);
\node[align=center,font=\sffamily\small\bfseries] at (0,2.45) {untrusted imports\\and tool output};
\node[align=center,font=\sffamily\small\bfseries] at (0,1.35) {validated\\context};
\node[align=center,font=\sffamily\small\bfseries] at (0,0) {trusted\\state};
\node[histbox, fill=white, draw=BitRose!70, minimum width=2.15cm, minimum height=.55cm, font=\sffamily\scriptsize] (inj) at (4.0,1.6) {injection};
\node[histbox, fill=white, draw=BitRose!70, minimum width=2.15cm, minimum height=.55cm, font=\sffamily\scriptsize] (stale) at (4.15,0.2) {stale state};
\node[histbox, fill=white, draw=BitRose!70, minimum width=2.15cm, minimum height=.55cm, font=\sffamily\scriptsize] (falsec) at (4.05,-1.35) {false completion};
\draw[figdanger] (inj.west) -- (2.0,1.15);
\draw[figdanger] (stale.west) -- (1.3,0.0);
\draw[figdanger] (falsec.west) -- (.85,-.55);
\end{tikzpicture}
\caption{Trust rises only after validation. Raw repository text is not automatically safe context.}
\label{fig:threat-surface}
\end{figure}


Continuity moves useful state into plain files. That improves visibility, but it also creates an obvious attack surface.

\subsection{Assets}

The protected assets are source code, secrets, project state, task queues, branch and review records, generated artifacts, and tool permissions.

\subsection{Adversaries}

Relevant adversaries include malicious pull requests, poisoned issue bodies, compromised documentation, prompt-injected tool output, careless agents, stale generated summaries, and over-permissioned automation.

\subsection{Failure modes}

The core failures are prompt injection, silent instruction override, stale-state obedience, false completion, secret persistence, and agent collision. Multi-agent use adds lock poisoning, fake claims, and adversarial handoff records.

\subsection{Mitigations}


\begin{table}[H]
\centering
\caption{Condensed threat model.}
\label{tab:threat-model}
\small
\rowcolors{2}{BitGray}{white}
\begin{tabularx}{\linewidth}{>{\raggedright\arraybackslash}p{2.75cm}X}
\rowcolor{BitNavy!8}
\textbf{Element} & \textbf{Examples} \\
Assets & source code, secrets, state files, task queues, artifacts, permissions \\
Adversaries & malicious pull requests, poisoned issues, prompt-injected docs, careless agents \\
Attack surfaces & manifests, state files, imported text, tool output, generated summaries, hooks \\
Failures & injection, stale-state obedience, false completion, secret persistence, agent collision \\
Mitigations & validation, context budgets, hooks, sandboxing, least privilege, review, signed commits \\
\end{tabularx}
\end{table}


The mitigation stack is straightforward in principle: treat imported text as untrusted, validate every trusted file, require evidence for completion, keep permissions narrow, and insist on review before high-impact state changes become canonical.
